\documentclass[11pt]{article}
\usepackage[utf8]{inputenc}
\usepackage[spanish, es-tabla]{babel}
\usepackage{amsmath, amssymb}
\usepackage{booktabs}
\usepackage{enumitem}
\usepackage[margin=2.5cm]{geometry}

\setlength{\parindent}{0pt}
\setlength{\parskip}{0.65em}

\title{Econometría ICS-294\\Guía de ejercicios: Clase 3}
\author{Francisco Alfaro Medina}
\date{}

\begin{document}
\maketitle

\section*{Objetivo}

Esta guía busca practicar los conceptos centrales de causalidad vistos en la clase: outcomes potenciales, contrafactuales, ATE, ATT, sesgo de selección y aleatorización.

\section*{Ejercicios}

\begin{enumerate}[label=\textbf{Ejercicio \arabic*.}, leftmargin=*]

\item \textbf{Asociación versus causalidad.}

Para cada afirmación, indique si describe una asociación o una interpretación causal. Justifique brevemente.
\begin{enumerate}[label=\alph*)]
    \item Las personas con más años de educación tienen, en promedio, salarios más altos.
    \item Un año adicional de educación aumenta el salario mensual promedio en \$80.000.
    \item Las casas con alarma reportan más robos que las casas sin alarma.
    \item Participar en un programa de capacitación aumenta la probabilidad de encontrar empleo.
\end{enumerate}

\item \textbf{Outcomes potenciales.}

Considere un programa de tutorías para estudiantes de primer año. Defina:
\begin{enumerate}[label=\alph*)]
    \item La variable de tratamiento $T_i$.
    \item El outcome observado $Y_i$.
    \item Los outcomes potenciales $Y_i^T$ y $Y_i^{NT}$.
    \item El efecto causal individual.
\end{enumerate}

Explique por qué el efecto causal individual no puede observarse directamente.

\item \textbf{Datos observados y datos contrafactuales.}

Complete la tabla indicando qué outcome potencial es observado y cuál es contrafactual.

\begin{center}
\begin{tabular}{lccc}
\toprule
Estudiante & $T_i$ & Puntaje observado $Y_i$ & ¿Qué no observamos? \\
\midrule
Ana & 1 & 78 & \\
Bruno & 0 & 65 & \\
Camila & 1 & 84 & \\
Diego & 0 & 70 & \\
Elena & 1 & 90 & \\
\bottomrule
\end{tabular}
\end{center}

\item \textbf{ATE y ATT.}

Use la notación de outcomes potenciales para escribir:
\begin{enumerate}[label=\alph*)]
    \item El efecto promedio del tratamiento, $ATE$.
    \item El efecto promedio del tratamiento sobre los tratados, $ATT$.
    \item La diferencia conceptual entre ambos.
\end{enumerate}

Luego explique por qué el $ATT$ contiene un término no observable.

\item \textbf{Comparación simple de promedios.}

Suponga que queremos estimar el efecto de asistir a la universidad sobre el salario mensual. En una muestra se observa:

\begin{center}
\begin{tabular}{lc}
\toprule
Grupo & Salario promedio mensual \\
\midrule
Asistió a la universidad & \$1.200.000 \\
No asistió a la universidad & \$850.000 \\
\bottomrule
\end{tabular}
\end{center}

\begin{enumerate}[label=\alph*)]
    \item Calcule la diferencia de promedios.
    \item ¿Bajo qué condición esta diferencia puede interpretarse como efecto causal?
    \item Mencione dos razones por las cuales esta comparación podría estar sesgada.
\end{enumerate}

\item \textbf{Descomposición del sesgo de selección.}

Parta de:
\[
\beta_1 = E(Y_i \mid T_i=1)-E(Y_i \mid T_i=0).
\]

Usando outcomes potenciales, muestre que:
\[
\beta_1 =
E(Y_i^T-Y_i^{NT}\mid T_i=1)
+
\left[
E(Y_i^{NT}\mid T_i=1)-E(Y_i^{NT}\mid T_i=0)
\right].
\]

Identifique cuál término corresponde al $ATT$ y cuál corresponde al sesgo de selección.

\item \textbf{Signo del sesgo.}

Para cada caso, indique si esperaría un sesgo de selección positivo, negativo o ambiguo. Justifique.
\begin{enumerate}[label=\alph*)]
    \item Efecto de asistir a la universidad sobre salarios.
    \item Efecto de instalar alarmas sobre robos en casas.
    \item Efecto de participar voluntariamente en un gimnasio sobre salud.
    \item Efecto de recibir clases de reforzamiento sobre notas finales.
\end{enumerate}

\item \textbf{Caso aplicado: alarmas y robos.}

Un municipio quiere evaluar si instalar alarmas reduce los robos domiciliarios. Se observa que las casas con alarma tienen más robos que las casas sin alarma.

\begin{enumerate}[label=\alph*)]
    \item ¿Significa esto necesariamente que las alarmas aumentan los robos?
    \item ¿Qué características podrían diferenciar a las casas que instalan alarmas?
    \item Explique cómo aparece el sesgo de selección en este ejemplo.
    \item Proponga una forma ideal de evaluar causalmente el efecto de las alarmas.
\end{enumerate}

\item \textbf{Aleatorización.}

Explique por qué asignar aleatoriamente un tratamiento permite interpretar la diferencia de promedios como efecto causal.

En su respuesta, use la condición:
\[
E(Y_i^{NT}\mid T_i=1)=E(Y_i^{NT}\mid T_i=0).
\]

\item \textbf{Experimento aleatorio.}

Una universidad tiene 400 estudiantes elegibles para recibir tutorías. Por cupos, solo 200 pueden recibirlas. Se asigna el tratamiento aleatoriamente.

\begin{enumerate}[label=\alph*)]
    \item ¿Cuál es el grupo tratado?
    \item ¿Cuál es el grupo de control?
    \item ¿Qué outcome podría usarse para evaluar el programa?
    \item ¿Qué comparación permitiría estimar el efecto causal promedio?
    \item ¿Por qué este diseño es mejor que dejar que los estudiantes se inscriban voluntariamente?
\end{enumerate}

\item \textbf{Lectura crítica de resultados.}

Un estudio reporta la siguiente regresión:
\[
\widehat{salario}_i = 420000 + 95000\, universidad_i,
\]
donde $universidad_i=1$ si la persona asistió a la universidad y $0$ si no.

\begin{enumerate}[label=\alph*)]
    \item Interprete el coeficiente $95000$ como diferencia condicional de medias.
    \item ¿Qué información adicional necesita para interpretarlo causalmente?
    \item Escriba una frase cuidadosa que no afirme causalidad.
    \item Escriba una frase causal que solo sería válida bajo supuestos fuertes.
\end{enumerate}

\item \textbf{Diseño de evaluación.}

Elija una política pública o programa de su interés. Puede ser una beca, subsidio, capacitación, campaña de salud, apoyo escolar u otro.

Describa:
\begin{enumerate}[label=\alph*)]
    \item El tratamiento.
    \item El outcome de interés.
    \item La población objetivo.
    \item El principal problema de selección que enfrentaría una evaluación observacional.
    \item Una estrategia experimental o cuasi experimental que podría ayudar a estimar un efecto causal.
\end{enumerate}

\end{enumerate}

\section*{Preguntas de cierre}

\begin{enumerate}[label=\arabic*.]
    \item ¿Cuál es la diferencia entre un contrafactual y un grupo de control?
    \item ¿Por qué comparar tratados y no tratados puede ser problemático?
    \item ¿Qué problema resuelve la aleatorización?
    \item ¿Por qué el $ATE$ y el $ATT$ no siempre coinciden?
\end{enumerate}

\end{document}
